\documentclass[11pt]{article}
\usepackage[utf8]{inputenc}
\usepackage{amsmath,amssymb}
\usepackage{hyperref}
\title{Scalable Pangenome Graph Reference for Large-Scale Settings}
\author{Anonymous}
\date{}
\begin{document}
\maketitle

\begin{abstract}
This paper studies Pangenome Graph Reference. Grounded in a real, citable literature \cite{ref1,ref2,ref3,ref4,ref5,ref6,ref7,ref8},
we propose a focused method and report \textbf{illustrative} experiments.
All references below are real and verifiable (OpenAlex/arXiv); the reported
numbers are simulated to illustrate intended behaviour.
\end{abstract}

\section{Introduction}
Pangenome Graph Reference is an active research area. This work targets a concrete gap identified from
the recent literature and contributes a method together with an evaluation protocol.

\section{Related Work (real literature)}
The following works, retrieved from public bibliographic APIs, frame our study:
\begin{itemize}
\item Martin et al. (2022) — "Ensembl 2023", Nucleic Acids Research (cited 1109).
\item Wang et al. (2022) — "The Human Pangenome Project: a global resource to map genomic diversity", Nature (cited 571).
\item Zhou et al. (2022) — "Graph pangenome captures missing heritability and empowers tomato breeding", Nature (cited 537).
\item Li et al. (2020) — "The design and construction of reference pangenome graphs with minigraph", Genome biology (cited 508).
\item Sirén et al. (2021) — "Pangenomics enables genotyping of known structural variants in 5202 diverse genomes", Science (cited 423).
\item Hickey et al. (2023) — "Pangenome graph construction from genome alignments with Minigraph-Cactus", Nature Biotechnology (cited 322).
\end{itemize}

\section{Method}
We decompose the problem across shards with a communication-efficient aggregation rule that keeps overhead sublinear in scale. We instantiate this idea for Pangenome Graph Reference and analyse its cost/quality trade-off.

\section{Experiments (Illustrative)}
\begin{center}
\begin{tabular}{lcc}
System & Primary Metric & Efficiency \\ \hline
Strong baseline & 0.812 & 1.00$\times$ \\
Ablation & 0.828 & 0.92$\times$ \\
\textbf{Ours} & \textbf{0.851} & \textbf{0.71$\times$}
\end{tabular}
\end{center}
\emph{Metrics simulated to illustrate the intended effect; not measured.}

\section{Conclusion}
We connected a real literature to a concrete method for Pangenome Graph Reference. Future work will
validate the illustrative results with real experiments.

\begin{thebibliography}{9}
\bibitem{ref1} Fergal J. Martin, M Ridwan Amode, Alisha Aneja et al.. Ensembl 2023. \emph{Nucleic Acids Research}, 2022. DOI:10.1093/nar/gkac958
\bibitem{ref2} Ting Wang, Lucinda Antonacci-Fulton, Kerstin Howe et al.. The Human Pangenome Project: a global resource to map genomic diversity. \emph{Nature}, 2022. DOI:10.1038/s41586-022-04601-8
\bibitem{ref3} Yao Zhou, Zhiyang Zhang, Zhigui Bao et al.. Graph pangenome captures missing heritability and empowers tomato breeding. \emph{Nature}, 2022. DOI:10.1038/s41586-022-04808-9
\bibitem{ref4} Heng Li, Xiaowen Feng, Chong Chu. The design and construction of reference pangenome graphs with minigraph. \emph{Genome biology}, 2020. DOI:10.1186/s13059-020-02168-z
\bibitem{ref5} Jouni Sirén, Jean Monlong, Xian Chang et al.. Pangenomics enables genotyping of known structural variants in 5202 diverse genomes. \emph{Science}, 2021. DOI:10.1126/science.abg8871
\bibitem{ref6} Glenn Hickey, Jean Monlong, Jana Ebler et al.. Pangenome graph construction from genome alignments with Minigraph-Cactus. \emph{Nature Biotechnology}, 2023. DOI:10.1038/s41587-023-01793-w
\bibitem{ref7} Erik Garrison, Zev Kronenberg, Eric T. Dawson et al.. A spectrum of free software tools for processing the VCF variant call format: vcflib, bio-vcf, cyvcf2, hts-nim and slivar. \emph{PLoS Computational Biology}, 2022. DOI:10.1371/journal.pcbi.1009123
\bibitem{ref8} Erich D. Jarvis, Giulio Formenti, Arang Rhie et al.. Semi-automated assembly of high-quality diploid human reference genomes. \emph{Nature}, 2022. DOI:10.1038/s41586-022-05325-5
\end{thebibliography}
\end{document}
